\chapter{Swollen or Painful Testicle}

\section*{Definition}
Enlargement, pain, or tenderness of one or both testicles, ranging from mild discomfort to acute severe pain requiring emergency evaluation.

\section*{What Happens in Your Body}
Testicular torsion (spermatic cord twisting) compromises blood flow causing reduced blood flow within 4-6 hours — surgical emergency. Epididymitis results from infection ascending from urethra/bladder. Hydrocele, varicocele, and spermatocele cause painless swelling. Testicular tumors typically present as painless solid mass.

\section*{Causes}
\begin{itemize}
  \item Testicular torsion (surgical emergency)
  \item Epididymitis (often bacterial or sexually transmitted)
  \item Orchitis (viral — mumps, or bacterial)
  \item Hydrocele (fluid collection around testicle)
  \item Varicocele (dilated scrotal veins)
  \item Spermatocele (epididymal cyst)
  \item Testicular tumor or cancer
  \item Testicular trauma
  \item Inguinal hernia (scrotal extension)
  \item Torsion of appendix testis (small testicular appendage)
\end{itemize}

\section*{When to See a Doctor}
See your doctor if:
\begin{itemize}
  \item Symptoms are severe or getting worse over time
  \item Sudden severe testicular pain with nausea/vomiting (possible torsion — immediate emergency), or painless testicular lump
  \item You have other concerning symptoms such as fever, weight loss, or bleeding
\end{itemize}

\section*{Related Symptoms}
\begin{itemize}
  \item Groin pain
  \item Scrotal swelling
  \item Nausea and vomiting
  \item Fever
  \item Urinary symptoms
\end{itemize}
\textbf{Important:} If this symptom persists, your doctor will check for serious underlying causes.

\section*{Self-Care / Home Management}
\begin{itemize}
  \item For mild scrotal discomfort without signs of torsion, NSAIDs (ibuprofen 400--600 mg three times daily) and scrotal elevation (supportive underwear or rolled towel) may provide relief.
  \item Apply an ice pack wrapped in cloth to the scrotum for 15 minutes at a time to reduce swelling in suspected epididymitis or trauma.
  \item Do not apply heat or attempt to manipulate a painful testicle; any acute testicular pain warrants immediate medical evaluation.
  \item If a painless testicular lump is discovered, no scrotal support or OTC treatment is appropriate; schedule medical evaluation promptly.
\end{itemize}

\section*{How Your Doctor Will Evaluate You}
\begin{itemize}
  \item Emergency evaluation for acute severe pain: assess the cremasteric reflex (absent in torsion), testicular lie (transverse or high-riding in torsion), and Prehn sign (relief with elevation suggests epididymitis).
  \item Scrotal Doppler ultrasound with colour flow is the definitive imaging study: detects absent or reduced flow in torsion, increased flow in epididymitis, or solid versus cystic nature of masses.
  \item Urinalysis and urine culture if UTI or epididymitis is suspected; urethral swab or urine NAAT for chlamydia and gonorrhoea in sexually active men.
  \item Testicular tumour markers (AFP, beta-hCG, LDH) if ultrasound demonstrates a solid intratesticular mass; radical inguinal orchiectomy is diagnostic and therapeutic.
\end{itemize}

\section*{Treatment Options}
\begin{itemize}
  \item Testicular torsion: immediate manual detorsion (outward rotation of the affected testicle) followed by bilateral orchidopexy within four to six hours of symptom onset to preserve testicular viability.
  \item Epididymitis: antibiotics based on aetiology; ceftriaxone 500 mg IM single dose plus doxycycline 100 mg BID for 10--14 days for STI pathogens; levofloxacin 500 mg daily for enteric organisms.
  \item Testicular tumour: radical inguinal orchiectomy with staging CT of chest, abdomen, and pelvis; further management (surveillance, chemotherapy, RPLND) based on histology and staging.
  \item Hydrocele, varicocele, and spermatocele are managed conservatively if asymptomatic; surgical repair (hydrocelectomy, varicocelectomy) is indicated if symptomatic or for fertility concerns.
\end{itemize}

\section*{References}
\begin{thebibliography}{99}
\bibitem{swollen_or_painful_testicle:aua-testicular-pain} Jarow JP, et al. ``AUA guideline for the evaluation and management of scrotal pain.'' \emph{J Urol}. 2024;211(3):412-423.
\bibitem{swollen_or_painful_testicle:uptodate-testicular-pain} Osterberg EC, et al. Evaluation of acute scrotal pain in adults. In: UpToDate, Post TW (Ed), UpToDate, Waltham, MA; 2025.
\bibitem{swollen_or_painful_testicle:nejm-torsion} Ringdahl E, et al. ``Testicular torsion.'' \emph{N Engl J Med}. 2023;389(2):152-160.
\bibitem{swollen_or_painful_testicle:bmj-testicular-pain} Powell JL, et al. ``Acute scrotal pain.'' \emph{BMJ}. 2024;384:e076543.
\bibitem{swollen_or_painful_testicle:harrison-testicular} Jameson JL, et al. \emph{Harrison\'s Principles of Internal Medicine}. 21st ed. McGraw-Hill; 2022. Chapter 360: Testicular Disorders.
\end{thebibliography}
